\documentclass[11pt]{article}

\usepackage[T1]{fontenc}
\usepackage[utf8]{inputenc}
\usepackage[margin=0.85in]{geometry}
\usepackage{amsmath,amssymb}
\usepackage{booktabs}
\usepackage{enumitem}
\usepackage{needspace}
\usepackage{xcolor}
\usepackage{listings}
\usepackage[colorlinks=true,linkcolor=blue,urlcolor=blue]{hyperref}
\usepackage{microtype}

\setlength{\parindent}{0pt}
\setlength{\parskip}{0.55em}
\setlist{nosep,leftmargin=*}
\emergencystretch=2em

\definecolor{codeblue}{RGB}{0,92,148}
\definecolor{codegreen}{RGB}{0,112,76}
\definecolor{codegray}{RGB}{90,98,108}
\definecolor{codewash}{RGB}{246,248,250}
\definecolor{outputwash}{RGB}{251,247,238}
\definecolor{rulegray}{RGB}{205,211,218}

\lstdefinestyle{mathematica-input}{
  basicstyle=\ttfamily\small,
  backgroundcolor=\color{codewash},
  frame=single,
  rulecolor=\color{rulegray},
  framerule=0.4pt,
  framesep=6pt,
  xleftmargin=0.6em,
  xrightmargin=0.6em,
  breaklines=true,
  breakatwhitespace=false,
  columns=fullflexible,
  keepspaces=true,
  showstringspaces=false,
  keywordstyle=\color{codeblue}\bfseries,
  commentstyle=\color{codegreen}\itshape,
  stringstyle=\color{codegray},
  morekeywords={
    AssociationMap,ConstantArray,FullSimplify,GogAvailableSizes,
    GogDiffusionState,GogRandomTrajectory,GogSmallModel,
    GogTransitionCounts,GogTransitionMatrix,GogValidateSmallModels,
    IdentityMatrix,MatrixForm,Needs,NullSpace,PacletDirectoryLoad,
    ReplacePart,SymmetricMatrixQ,Table,Transpose
  }
}

\lstdefinestyle{mathematica-output}{
  basicstyle=\ttfamily\small,
  backgroundcolor=\color{outputwash},
  frame=single,
  rulecolor=\color{rulegray},
  framerule=0.4pt,
  framesep=6pt,
  xleftmargin=0.6em,
  xrightmargin=0.6em,
  breaklines=true,
  breakatwhitespace=false,
  columns=fullflexible,
  keepspaces=true,
  showstringspaces=false
}

\newcommand{\one}{\mathbf 1}
\newcommand{\V}{\mathcal V}
\newcommand{\GoG}{\Gamma}
\newcommand{\mathematicacells}[1]{%
  \par\medskip
  \noindent\makebox[\linewidth][c]{\small Mathematica cells: #1}\par
  \vspace{-0.3em}%
}

\title{Whitehead Field Dynamics\\
  \large Computational Notebook Notes and Exact Mathematica Cells}
\author{Alexander Kolpakov}
\date{September 21, 2026}

\begin{document}
\maketitle

\begin{abstract}
This companion note turns the executable parts of
\path{whitehead-field-dynamics.nb} into a short, reproducible transcript.  It
records the local pair-averaging calculation, the exact small graph-of-graphs
models used by the notebook, the two transition clocks, an exact diffusion
snapshot, one seeded sample path, and the finite consensus products for
$K_4$ and $Q_3$.  Every displayed number is obtained with exact rational
arithmetic before a decimal approximation is requested.  The note is not a
replacement for the full manuscript: its purpose is to place selected
Mathematica cells and their interpretation in one compilable \TeX{} file.
\end{abstract}

\section{What is being computed}

The state space consists of isomorphism classes of connected cubic
multigraphs on an even number $N$ of vertices; loops and parallel edges are
allowed.  A Whitehead move selects a carrier edge, collapses it, and chooses
one of the two other pairings of the four exposed half-edges.  A scalar field
$f\in\mathbb R^N$ is carried along with the graph.  If the selected edge has
endpoints $i,j$, the field update is
\[
  (f_i,f_j)\longmapsto
  \left(\frac{f_i+f_j}{2},\frac{f_i+f_j}{2}\right).
\]

Three computational objects must be kept separate.

\begin{enumerate}
  \item The \emph{carrier field} uses a concrete edge and, after an
  unlabelled move, a settlement permutation for its coordinates.
  \item Rule A is the quotient chain that retains all $3N$ compressed
  resolution slots, including multiplicity and holding slots.
  \item Fixed-clock Rule B assigns one of the same $3N$ ticks to each
  distinct nonself quotient neighbour and turns every unused tick into a
  hold.  Its jump chain, denoted \texttt{"BJump"} in the code, is Rule B
  conditioned on moving and is a different Markov chain.
\end{enumerate}

The notebook's interactive laboratory bundles complete quotient data only for
$N=2,4,6$.  The larger finite tables in the manuscript come from the Python
enumeration pipeline and should not be confused with the compact notebook
resource.

\section{A visual map of the notebook}

The refreshed notebook is deliberately a companion rather than a second copy
of the paper.  Its seven sections can be read in sequence, but each visual is
also a useful entry point into one computational question.

\begin{center}
\small
\begin{tabular}{@{}p{0.27\linewidth}p{0.66\linewidth}@{}}
\toprule
Notebook visual & Computational role \\
\midrule
Visual abstract & Separates the carrier field, the fixed-labelled move graph,
  the multiplicity chain, and the simple-support chain. \\
Whitehead move & Shows the selected carrier edge, its four exposed half-edges,
  and the two compressed resolutions. \\
Variance drop & Turns one pair average into the exact loss
  $(f_i-f_j)^2/2$. \\
Consensus schedules & Displays the matching layers that give finite consensus
  on $K_4$ and $Q_3$. \\
Transition heat maps & Compares Rule A with fixed-clock Rule B on the same
  five quotient classes at $N=4$. \\
Diffusion laboratory & Places exact rational mass and one independently
  sampled path on the quotient graph. \\
Enumeration growth & Compares class counts and Rule-A slot records through
  $N=14$ on a logarithmic scale. \\
\bottomrule
\end{tabular}
\end{center}

The notebook keeps explanatory prose short and places executable cells beside
the relevant picture.  In this note, every input/output pair has one shared
caption of the form ``Mathematica cells: \emph{topic}.''  Input cells retain
the pale-blue background used in the notebook, while exact outputs retain the
warm cream background.  This makes the direction of a calculation visible
without repeating labels such as ``input'' and ``output'' above every box.

\section{The local averaging calculation}

For a non-loop pair $\{i,j\}$, the averaging matrix is
\[
  A_{ij}=I-\frac12(e_i-e_j)(e_i-e_j)^T.
\]
It is a symmetric idempotent matrix, fixes $\one$, and hence is the orthogonal
projection onto $\{f:f_i=f_j\}$.  The two-coordinate energy calculation can
be checked without constructing the whole matrix.

\mathematicacells{one-pair loss}
\begin{lstlisting}[style=mathematica-input]
Clear[x, y];
FullSimplify[
  x^2 + y^2 - 2 ((x + y)/2)^2,
  Assumptions -> Element[{x, y}, Reals]
]
\end{lstlisting}

\begin{lstlisting}[style=mathematica-output]
(x - y)^2/2
\end{lstlisting}

Since the sum of the two coordinates is unchanged, the mean is unchanged as
well.  Applying the identity to the centered field gives the exact Lyapunov
law
\[
  \V(A_{ij}f)=\V(f)-\frac12(f_i-f_j)^2,
  \qquad
  \V(f)=\|f-\bar f\one\|_2^2.
\]
The following definition is used again for the consensus calculation.

\mathematicacells{averaging matrix}
\begin{lstlisting}[style=mathematica-input]
pairAverage[n_, {i_, j_}] := ReplacePart[
  IdentityMatrix[n],
  {{i, i} -> 1/2, {i, j} -> 1/2,
   {j, i} -> 1/2, {j, j} -> 1/2}
];

a = pairAverage[4, {1, 3}];
{
  SymmetricMatrixQ[a],
  a . a === a,
  a . ConstantArray[1, 4] === ConstantArray[1, 4]
}
\end{lstlisting}

\begin{lstlisting}[style=mathematica-output]
{True, True, True}
\end{lstlisting}

A settlement after a quotient move is a permutation matrix $P_\sigma$.
Because $P_\sigma$ preserves sums and Euclidean norms, replacing $A_{ij}$ by
$P_\sigma A_{ij}$ does not change this one-step variance identity.  It does,
however, affect which persistent coordinate pair is averaged later, so a
quotient trajectory alone is not enough to reconstruct the full field
cocycle.

\Needspace{22\baselineskip}
\section{Loading and validating the notebook models}

From a notebook opened at the repository root, the local companion paclet can
be loaded as follows.

\mathematicacells{local paclet}
\begin{lstlisting}[style=mathematica-input]
repo = NotebookDirectory[];
PacletDirectoryLoad[
  FileNameJoin[{repo, "Wolfram", "GogFieldDynamics"}]
];
Needs["GogFieldDynamics`"];

GogAvailableSizes[]
\end{lstlisting}

\begin{lstlisting}[style=mathematica-output]
{2, 4, 6}
\end{lstlisting}

For each bundled size, the next query records
\[
  \{N,\ \text{number of classes},\ \text{support edges},\
      \text{total Rule-A slots}\}.
\]

\mathematicacells{resource summary}
\begin{lstlisting}[style=mathematica-input]
Table[
  With[{m = GogSmallModel[n]},
    {n, m["NodeCount"], m["EdgeCount"],
     m["RuleASlotCount"]}
  ],
  {n, GogAvailableSizes[]}
]
\end{lstlisting}

\begin{lstlisting}[style=mathematica-output]
{{2, 2, 1, 12}, {4, 5, 4, 60}, {6, 17, 25, 306}}
\end{lstlisting}

The class counts $2,5,17$ agree with the corresponding first terms of OEIS
A005967.  The last column also passes the independent slot check
\[
  \#\text{classes}\times 3N:\qquad
  2\cdot6=12,\quad5\cdot12=60,\quad17\cdot18=306.
\]
The paclet validator checks completeness metadata, node and slot counts,
support edges, row sums, stochasticity, Rule-B symmetry, and uniform Rule-B
stationarity.

\mathematicacells{compact validation report}
\begin{lstlisting}[style=mathematica-input]
r = GogValidateSmallModels[];
{
  r["Valid"],
  AssociationMap[
    r["Models", #]["Valid"] &,
    GogAvailableSizes[]
  ]
}
\end{lstlisting}

\begin{lstlisting}[style=mathematica-output]
{True, <|2 -> True, 4 -> True, 6 -> True|>}
\end{lstlisting}

These checks validate the bundled finite data and the implemented matrix
conventions.  They are not an asymptotic theorem.

\section{An exact five-state example}

At $N=4$ there are five unlabelled carrier classes.  Rule A retains the
integer multiplicity of every compressed resolution.  Its count matrix is

\mathematicacells{Rule-A counts}
\begin{lstlisting}[style=mathematica-input]
countsA = GogTransitionCounts[4, "A"];
{countsA, Total /@ countsA}
\end{lstlisting}

\begin{lstlisting}[style=mathematica-output]
{{{4, 4, 4, 0, 0},
  {6, 6, 0, 0, 0},
  {2, 0, 8, 0, 2},
  {0, 0, 0, 6, 6},
  {0, 0, 4, 2, 6}},
 {12, 12, 12, 12, 12}}
\end{lstlisting}

Dividing by $3N=12$ gives the exact transition matrix
\[
P_A=
\begin{pmatrix}
1/3&1/3&1/3&0&0\\
1/2&1/2&0&0&0\\
1/6&0&2/3&0&1/6\\
0&0&0&1/2&1/2\\
0&0&1/3&1/6&1/2
\end{pmatrix}.
\]
This matrix is not symmetric.  Solving $\pi_A P_A=\pi_A$ with exact
arithmetic gives

\mathematicacells{Rule-A stationary law}
\begin{lstlisting}[style=mathematica-input]
pA = GogTransitionMatrix[4, "A"];
piA = First[
  NullSpace[Transpose[pA] - IdentityMatrix[Length[pA]]]
];
piA/Total[piA]
\end{lstlisting}

\begin{lstlisting}[style=mathematica-output]
{1/5, 2/15, 2/5, 1/15, 1/5}
\end{lstlisting}

The nonuniform vector is the finite $N=4$ instance of the projected cubic
configuration-model law.  Treating Rule A as uniform on isomorphism classes
would therefore change the process.

\Needspace{22\baselineskip}
For fixed-clock Rule B, every distinct nonself support neighbour has
probability $1/12$ and the remaining ticks hold:

\mathematicacells{fixed-clock Rule B}
\begin{lstlisting}[style=mathematica-input]
pB = GogTransitionMatrix[4, "B"];
MatrixForm[pB]
\end{lstlisting}

\begin{lstlisting}[style=mathematica-output]
{{5/6, 1/12, 1/12, 0, 0},
 {1/12, 11/12, 0, 0, 0},
 {1/12, 0, 5/6, 0, 1/12},
 {0, 0, 0, 11/12, 1/12},
 {0, 0, 1/12, 1/12, 5/6}}
\end{lstlisting}

The defining structural checks are visible directly in the output:

\mathematicacells{symmetry and stationarity}
\begin{lstlisting}[style=mathematica-input]
u = ConstantArray[1/5, 5];
{SymmetricMatrixQ[pB], u . pB}
\end{lstlisting}

\begin{lstlisting}[style=mathematica-output]
{True, {1/5, 1/5, 1/5, 1/5, 1/5}}
\end{lstlisting}

Thus the fixed-clock chain is symmetric and the uniform class law is
stationary.  It also satisfies
\[
  I-P_B=\frac{1}{3N}L_{\mathrm{supp}}.
\]
Holding cannot simply be discarded.  Already at $N=2$, the fixed clock and
the chain conditioned on moving are

\mathematicacells{Rule B versus its jump chain}
\begin{lstlisting}[style=mathematica-input]
{
  GogTransitionMatrix[2, "B"],
  GogTransitionMatrix[2, "BJump"]
}
\end{lstlisting}

\begin{lstlisting}[style=mathematica-output]
{{{5/6, 1/6}, {1/6, 5/6}},
 {{0, 1}, {1, 0}}}
\end{lstlisting}

The two matrices have the same nonself support, but different clocks,
spectra, and path probabilities.

\section{Exact diffusion and a sampled path}

The blue node fill in the notebook laboratory is computed by exact row-vector
multiplication.  Its default static view uses $N=4$, Rule A, start class $1$,
and time $t=4$.

\mathematicacells{default exact snapshot}
\begin{lstlisting}[style=mathematica-input]
p4A = GogDiffusionState[4, "A", 1, 4];
{p4A, N[p4A, 6]}
\end{lstlisting}

\begin{lstlisting}[style=mathematica-output]
{{185/648, 49/216, 59/162, 1/54, 17/162},
 {0.285494, 0.226852, 0.364198, 0.0185185, 0.104938}}
\end{lstlisting}

All five classes are reachable by time four, but the mass is far from the
stationary vector.  Running the same initial condition and time with the
fixed-clock Rule B gives

\mathematicacells{Rule-B comparison}
\begin{lstlisting}[style=mathematica-input]
p4B = GogDiffusionState[4, "B", 1, 4];
{p4B, N[p4B, 6]}
\end{lstlisting}

\begin{lstlisting}[style=mathematica-output]
{{5623/10368, 1181/5184, 4121/20736,
  41/20736, 151/5184},
 {0.542342, 0.227816, 0.198736, 0.00197724, 0.0291281}}
\end{lstlisting}

The extra holding probability in Rule B is plainly visible: after four ticks,
more than half of the mass remains at the starting class.  This comparison is
about two quotient clocks; it does not yet specify a Rule-B lift that chooses
which concrete carrier edge to average.

The orange ring and trail in the notebook show one random realization, not an
estimate of the blue distribution.  With the notebook's default seed, the
native sampler returns

\mathematicacells{one reproducible trajectory}
\begin{lstlisting}[style=mathematica-input]
GogRandomTrajectory[4, "A", 1, 12, 2026]
\end{lstlisting}

\begin{lstlisting}[style=mathematica-output]
{1, 3, 1, 3, 3, 3, 3, 5, 3, 1, 2, 2, 1}
\end{lstlisting}

Repeated class labels are genuine holding events.  Rule-A multiplicity is
retained because the sampler chooses from the exact transition matrix (and
the optional SetReplace engine repeats rewrite rules according to the same
integer count matrix).

\Needspace{30\baselineskip}
\section{Finite consensus words}

Pair averages on disjoint edges commute, so a perfect matching can be treated
as one block.  The notebook visualizes two such blocks on $K_4$ and three
coordinate matchings on the cube $Q_3$.  The following exact calculation
checks that the corresponding products are the consensus matrices
$J_4=\one\one^T/4$ and $J_8=\one\one^T/8$.

\mathematicacells{matching products}
\begin{lstlisting}[style=mathematica-input]
matching[n_, edges_] :=
  Dot @@ (pairAverage[n, #] & /@ edges);

k4Word =
  matching[4, {{1, 3}, {2, 4}}] .
  matching[4, {{1, 2}, {3, 4}}];

q3Word =
  matching[8, {{1, 5}, {2, 6}, {3, 7}, {4, 8}}] .
  matching[8, {{1, 3}, {2, 4}, {5, 7}, {6, 8}}] .
  matching[8, {{1, 2}, {3, 4}, {5, 6}, {7, 8}}];

{
  k4Word === ConstantArray[1/4, {4, 4}],
  q3Word === ConstantArray[1/8, {8, 8}]
}
\end{lstlisting}

\begin{lstlisting}[style=mathematica-output]
{True, True}
\end{lstlisting}

For example, the first product is returned exactly as

\mathematicacells{the $K_4$ product}
\begin{lstlisting}[style=mathematica-input]
k4Word // MatrixForm
\end{lstlisting}

\begin{lstlisting}[style=mathematica-output]
{{1/4, 1/4, 1/4, 1/4},
 {1/4, 1/4, 1/4, 1/4},
 {1/4, 1/4, 1/4, 1/4},
 {1/4, 1/4, 1/4, 1/4}}
\end{lstlisting}

Consequently every initial field reaches its exact mean after the displayed
schedule.  This is stronger than exponential convergence, and it explains
the exceptional finite-time consensus at $N=4$ and $N=8$.  It should not be
extrapolated to arbitrary $N$: pair averaging has a dyadic denominator
obstruction, and a three-coordinate field such as $(0,0,1)$ cannot reach its
mean exactly after a finite word of pair averages.

\section{Interpretation limits}

The transcript supports several precise conclusions:

\begin{itemize}
  \item pair averaging is an orthogonal projection with exact variance loss;
  \item each Rule-A row contains exactly $3N$ resolution slots;
  \item the $N=4$ Rule-A class law is nonuniform;
  \item fixed-clock Rule B is symmetric and has uniform class stationarity;
  \item the interactive diffusion is exact rational matrix evolution; and
  \item the $K_4$ and $Q_3$ schedules are exact finite consensus words.
\end{itemize}

There are also three important non-conclusions.

\begin{enumerate}
  \item The bundled quotient model omits selected carrier edges, resolution
  identifiers, and settlement permutations.  It therefore cannot by itself
  produce a pathwise transported field or its Lyapunov exponent.
  \item A single sampled trajectory does not approximate the exact
  distribution.  The notebook deliberately renders them as separate layers.
  \item Complete finite checks, whether through $N=6$ in the notebook or
  through $N=14$ in the enumeration scripts, do not establish an asymptotic
  spectral or mixing statement.
\end{enumerate}

These boundaries are useful when reading the graphics: the notebook is an
exact laboratory for the specified finite quotient chains, not a substitute
for the manuscript's proofs or for a fully specified field lift.

\Needspace{18\baselineskip}
\section{Reproduction}

The displayed transcript was regenerated on September 21, 2026 with Wolfram
15.0.0 for macOS ARM (kernel build dated May 19, 2026).

The deterministic command-line transcript is produced by
\path{scripts/notebook_report_outputs.wls}:

\begin{lstlisting}[style=mathematica-input,title={Reproduction command: cell transcript}]
/Applications/Wolfram.app/Contents/MacOS/WolframKernel \
  -noinit -noprompt \
  -script scripts/notebook_report_outputs.wls
\end{lstlisting}

The compact data resource can be regenerated and byte-checked without a
Wolfram installation:

\begin{lstlisting}[style=mathematica-input,title={Reproduction command: data check}]
python3 scripts/build_small_gog_models.py --check
\end{lstlisting}

Compile this note from the repository root with

\begin{lstlisting}[style=mathematica-input,title={Reproduction command: TeX build}]
latexmk -pdf -interaction=nonstopmode -halt-on-error \
  whitehead-field-dynamics-computational-notes.tex
\end{lstlisting}

Build and test the executable mini-notebook version with

\begin{lstlisting}[style=mathematica-input,title={Reproduction command: executable mini-notebook}]
make computational-notes-notebook
make computational-notes-notebook-check
\end{lstlisting}

Then open {\small\path{whitehead-field-dynamics-computational-notes.nb}} and choose
\texttt{Evaluation > Evaluate Notebook}.  A complete successful run ends with
\texttt{AllTestsPassed -> True} in the final Mathematica output cell.

For the compact executable companion, its seven native graphics, and the
playable diffusion panel, open \path{whitehead-field-dynamics.nb}.

Mathematical proofs remain in \path{whitehead-field-dynamics.tex}.  The
complete claim-status distinctions remain in \path{docs/status.rst}.

\begingroup
\small
\begin{thebibliography}{9}

\bibitem{gdp-kolpakov}
Laura Grave de Peralta and Alexander Kolpakov,
\emph{Expansion properties of Whitehead moves on cubic graphs},
Comptes Rendus Math\'ematique \textbf{362} (2024), 1825--1836.
\url{https://doi.org/10.5802/crmath.691}

\bibitem{li}
Michael Li,
\emph{Increasingly Many Bounded Eigenvalues of the Graph of Whitehead Moves},
arXiv:2405.14592.
\url{https://arxiv.org/abs/2405.14592}

\end{thebibliography}
\endgroup

\end{document}
